% Chapter Template

\chapter{Background} % Main chapter title

\label{Background} % Change X to a consecutive number; for referencing this chapter elsewhere, use \ref{ChapterX}


%-----------------------------------
%	SECTION 1
%-----------------------------------
\section{Introduction to Preposition Supersenses}
\subsubsection{*Adpositions signal semantic relations}
\subsubsection{*These relations are ambiguous}
\subsubsection{*Disambiguations could be beneficial to downstream tasks}
\subsection{The SNACS Scheme}
\subsubsection{SNACS Label Hierarchy}
\subsubsection{Construal Analysis}
\subsection{The STREUSLE corpus}
\subsection{Previous Approaches}
\subsection{PSS in This Work}
\pagebreak

%-----------------------------------
%	SECTION 2
%-----------------------------------
\section{Neural Networks for NLP}

\begin{figure}
  \centering
    \includegraphics[width=0.5\textwidth]{Electron}
  \caption{ 
    (i) A single neuron with N inputs. \\
    (ii) A feed-forward network with a single output neuron \\
    (ii) A multiclass network architecture: A feed-forward network with a Softmax output layer of $l$ neurons \\    
  }
  \label{fig:networks}
\end{figure}

Inspired by the neural networks of the biological brain, Artificial Neural Networks are a family of computational models widely used today across many different fields of Artificial Intelligence research. Neural-based machine learning models have consistently been able to achieve and outperform state-of-the-art results when obtained using previous approaches. In the past decade, they have been a key aspect of several breakthroughs in computer vision and speech recognition, which led to a growing interest and popularity among other research fields including Natural Language Processing. This section lays out the foundations required for understanding the use of neural networks throughout this work.

\medskip

A single neuron is the basic element of a neural network. Similarly to a biological neuron in the human brain, it emits a real-valued signal based on input signals it receives from other neurons it is connected to. The emitted signal could be fed into other neurons at an upper layer, and so on, thus forming a network. Input signals emitted into the bottom layer of neurons are considered the input of the network (image pixels, sound wave amplitudes, etc). Similarly, the outputs emitted from the top-most layer are considered the output of the network (class probabilities, regression value, etc). Typically, each neuron represents some detectable feature of the input. The strength of the emitted signal represents the certainty of the detected feature, or lack thereof. The hierarchical structure of the network allows low-level features, detected by neurons at lower layers, to be used for detection of more complex features by neurons at higher layers, and so on.

\medskip

\subsection{A Single Neuron}

The output of a single neuron (figure \ref{fig:networks} (i)) is a composition of a weighted sum of its inputs (and perhaps a bias term), followed by a non-linearity - the Activation Function. Equation \ref{eq:neuron} provides a formal definition of the neuron function:
\begin{equation}
y = \sigma(\vec{w}\cdot\vec{x} + b)
\label{eq:neuron}
\end{equation}

Where \(\vec{x}\) is a vector of inputs, \(\vec{w}\) is a vector of weights, \(b\) is a bias term and \(\sigma\) is a non-linearity. Informally speaking, the weights determine which of the input signals the neuron is more sensitive to, and the bias term is used to adjust the weighted sum of the signals prior to the non-linearity. The use of non linear activation functions allows the network to model real world data, which is often complex and non linear. Figure \ref{fig:activations} describes two commonly used activation functions: $tanh$, and Rectified Linear Unit ($ReLU$). $tanh$ maps the input into the range of $(-1,1)$, which is useful for neurons detecting a binary feature, or networks used for binary classification. A neuron with the $ReLU$ activation will output a positive signal once the weighted sum of its inputs meets a certain threshold determined by the bias term, and zero otherwise. In other words, a $ReLU$ neuron will only "fire" if its input signals are strong enough. 

\begin{figure}
  \centering
    \includegraphics[width=0.5\textwidth]{Electron}
  \caption{
  Commonly used activation functions \\
  (i) $ ReLU(z) = max(0, z) $\\ 
  (ii) $ tanh(z) = \frac{2}{1 + e^{-2x}} - 1 $ \\
  }
  \label{fig:activations}
\end{figure}

\subsection{Feed-forward networks}

A set of \(k\) neurons operating on the same vector of input signals and sharing the same activation function \(\{y_i = \sigma(\vec{w_i}\cdot\vec{x} + b_i)\}\), is considered a single layer. Following eq. \ref{eq:neuron}, such layer computes an affine transformation on a vector of \(n\) input signals  \(\vec{x}_{nx1}\), followed by an element-wise non-linearity, as described in equation \ref{eq:neuronlayer}. 

\begin{equation}
    \vec{y}_{kx1} = \vec{\sigma}({W_{kxn}\vec{x}_{nx1}} + \vec{b}_{kx1})
    \label{eq:neuronlayer}
\end{equation}

Where the \(i\)-th row of \(W_{kxn}\) and \(\vec{b}_{kx1}\) are respectively the weight vector and the bias term of the \(i\)-th neuron, and \(\vec{\sigma}\) performs element-wise computation of the activation function \(\sigma\). Layers of neurons can be stacked on top of each other by setting the output of one layer to be the input to the next, as demonstrated in eq. \ref{eq:feedforward}:
\begin{equation}
    \begin{split}
    \vec{y}^1_{kx1} = \vec{\sigma}^1({W^1_{kxn}\vec{x}_{nx1}} + \vec{b}^1_{kx1}) \\
    \vec{y}^2_{kx1} = \vec{\sigma}^2({W^2_{kxk}\vec{y}^1_{kx1}} + \vec{b}^2_{kx1}) \\
    \vec{y}^3_{kx1} = \vec{\sigma}^3({W^3_{kxk}\vec{y}^2_{kx1}} + \vec{b}^3_{kx1}) \\
    out = \vec{\sigma}^{out}({W^{out}_{1xk}\vec{y}^3_{kx1}} + b^{out})
    \end{split}
    \label{eq:feedforward}
\end{equation}

Eq. \ref{eq:feedforward} describes a network with an input layer of size \(n\) followed by three layers of size \(k\) and an output layer consisting of a single neuron. Each layer has its own set of weights, bias terms and activation function. Neural networks consisting of such architecture - i.e a sequence of neuron layers stacked on top of each other, are known as \textbf{Feed Forward Networks}, or \textbf{Multi Layered Perceptrons} (due to the similarity of the early Perceptron model and a neuron). In feed-forward networks, information flows in one direction starting from the input layer, through intermediate layers (a.k.a hidden layers) and the output layer. This is in contrast to Recurrent Neural Networks (described in section todo), where neuron connections could form a cycle. A feed-forward network with a structure similar to that which defined in eq. \ref{eq:feedforward} is shown in figure \ref{fig:networks} (ii).

\subsection{The Softmax Layer}
A Softmax layer is a neuron layer that uses the Softmax activation function, which squashes the outputs of the layer to form a probability distribution. It is mostly used in the common case of multiclass classification, where each sample is to be assigned with a label from a predefined set of labels. Assuming a label set of size $L$, a Softmax layer consisting of $L$ neurons will generate a probability estimation for each label in the set. In such scenario, this layer would be the output layer of the network.

The Softmax function is different from the previously mentioned activation functions, in the sense that it operates on an entire vector of signals rather than on each individual signal independently. It receives a vector of the (intermediate) outputs of a neuron layer as input, and it outputs a vector of probabilities of the same dimension. The Softmax activation function is defined as follows:

\begin{equation}
    Softmax( \vec{x} )_j = \frac{e^{\vec{x}_j}}{\sum_{i=1}^k e^{\vec{x}_i}}
    \label{eq:softmax}
\end{equation}

Where $x \in \mathbb{R}^k$ is a vector of input signals. The equation of a Softmax \emph{layer} is obtained by using Softmax as the activation function \(\vec{\sigma}\) of eq.  \ref{eq:neuronlayer}. A full outline of a multiclass network architecture using a Softmax layer is outlined in figure \ref{fig:networks} (iii).

\subsection{Training a Neural Network}
Given a neural network architecture, the process of training the network aims to find a good set of \emph{model parameters} for a particular task. The parameters of a network are defined to be the entire set of weights and biases of neurons in the network. For instance, the set of parameters of the network defined in eq. \ref{eq:feedforward} is $\theta = \{W^1, \vec{b^1}, W^2, \vec{b^2}, W^3, \vec{b^3}, W^{out}, \vec{b^{out}}\}$. We can therefore consider $\theta$ as a vector of real-valued parameters by ordering and flattening the matrices and vectors in the set. Then, a neural network can be thought of as a function of an input sample $\vec{x}$ and a vector of model parameters $\theta$: 
$$ \vec{y} = \textbf{f}(\vec{x}, \theta) $$ 
In the supervised setting, we have a set of $N$ training samples $T = \{(\vec{x_i}, \vec{y_i})\}$ and we wish to obtain the best parameter vector $\hat{\theta}$ such that $\textbf{f}(\vec{x_i}, \hat{\theta})$ would be a good predictor of $y_i$ on average across all samples $(\vec{x_i}, \vec{y_i}) \in T$.  We use a \emph{loss function} $\mathcal{L}$ to measure the inconsistency between the ground truth $\vec{y_i}$ and the prediction $\textbf{f}(\vec{x_i}, \theta)$, which is also known as the \emph{cost} or \emph{loss} of the prediction. The learning process aims to minimize the average loss on the entire training set\footnote{In some cases, a regularization term $R(\theta)$ is added to the loss function as a mean to prevent overfitting. This is discussed in section todo.}:

\begin{equation}
\hat{\theta} = \argmin_{\theta} \frac{1}{N}\sum_{i=1}^N \mathcal{L}(\textbf{f}(\vec{x_i}, \theta), y_i)
\label{eq:bestparams}    
\end{equation}

\subsubsection{Stochastic Gradient Descent with Backpropagation}
While computing an exact solution for eq \ref{eq:bestparams} is often intractable, there are many optimizers aiming to find a good estimation of the best parameter vector $\hat{\theta}$. Most of these optimizers are gradient-based, and are typically variations of the Gradient Descent (GD) minimization algorithm. GD is an iterative process aiming to minimize a differentiable function $t$. It begins with a randomly assigned $\theta$. In each iteration, the gradient $\nabla t(\theta)$ is computed, and $\theta$ is updated accordingly:

$$ \theta \leftarrow \theta - \textit{lr} \cdot \nabla t(\theta) $$

This update step pushes $\theta$ in the opposite direction of the gradient, thus aiming to decrease $t(\theta)$. This process is repeated until a local minima is reached, a state which is practically detected by observing that $t(\theta)$ has not significantly decreased in the last window of iterations. The hyperparameter \textit{lr} - the \textit{learning rate} - controls the magnitude of the update step. A small learning rate may result in slow convergence, while a large learning rate may prevent the process from converging at all (as the magnified gradient could push $\theta$ too far and actually \emph{increase} $t(\theta)$). In some variations of GD, the learning rate is multiplied by a decay factor after each predefined number of iterations. 

Following eq \ref{eq:bestparams}, the target function we wish to minimize in order to train the network is:

$$ t(\theta) = \frac{1}{N}\sum_{i=1}^N \mathcal{L}(\textbf{f}(\vec{x_i}, \theta), y_i) $$

For $t$ to be differentiable with regards to $\theta$, we require choosing a neural network $f$ and a loss function $\mathcal{L}$ such that $f$ is differentiable with regards to $\theta$, and $\mathcal{L}$ is differentiable with regards to its first argument - the predicted vector.

In practice, computing the gradient of $t$ is very computationally expensive for a large training set, as it requires computing the gradient of the loss of each training sample separately. Instead, in a process known as \textbf{Mini-Batch Gradient Descent} or \textbf{Stochastic Gradient Descent (SGD)}, an approximation of the gradient is computed. The set of input samples is divided into batches of size $BATCH\_SIZE$. At each step of the iteration, the gradient is computed on the average loss of only the next batch of samples:

$$ \hat{t}(\theta) = \frac{1}{BATCH\_SIZE}\sum_{(\vec{x_i}, \vec{y_i}) \in Batch} \mathcal{L}(\textbf{f}(\vec{x_i}, \theta), \vec{y_i})  $$

$$ \theta \leftarrow \theta - \textit{lr} \cdot \nabla \hat{t}(\theta) $$

This update step can be thought of as training the network on a batch of samples. An \emph{epoch} represents a full pass of training over the entire set of batches in the training set. It is very common to train a network through several epochs, while shuffling the order of samples between each epoch. The right amount of epochs depends on the convergence rate and the size of the training set.

The following are a few commonly used variations of SGD:
\begin{itemize}
    \item \textbf{Adaptive Gradient Algorithm (AdaGrad)}: ???todo???
    \item \textbf{Root Mean Square Propagation (RMSProp)}: ???todo???
    \item \textbf{Adaptive Moment estimation (Adam}: ???todo???    
\end{itemize}

SGD (along with its variations) is a heuristic process and is not guaranteed to converge to a global minima. Emprical results consistently show, however, that neural networks can be trained successfully even when using plain SGD.

The gradient for a single sample in the training set -  $\nabla_\theta \mathcal{L}(\textbf{f}(\vec{x_i}, \theta), \vec{y_i})$ - can be computed efficiently in a process known as \textbf{Backpropagation}. The Backpropagation
algorithm is able to efficiently compute the gradient of a real function $t(\theta)$, provided that it holds several key properties: (a) $t$ is a composition of real functions, (b) each function takes only elements from $\theta$ or an output of another function as inputs, and (c) the partial derivatives of each function can be computed efficiently.  

It follows from (a), (b) that $t$ can be described as a directed acyclic graph, where each graph node $\textbf{v}$ represents an element from $\theta$ (a \textit{$\theta$-node}), or an output of an intermediate function (where the function and output are notated as $\textbf{f}_{\textbf{v}}$ and $o_{\textbf{v}}$ respectively). An edge $\textbf{u}\rightarrow\textbf{v}$ is present when $o_{\textbf{u}}$ is an input of $\textbf{f}_\textbf{v}$. It holds for the (single) sink node $\textbf{s}$ that $o_{\textbf{s}} = t(\theta)$. The \textit{$\theta$-nodes} are the only ones with no incoming edges. 

The Backpropagation process computes a partial derivative value of $t$ for each \emph{edge} in the graph. It begins with edges connected to the sink node, and going backwards in reversed topological order (utilizing the chain rule for function derivation). The partial derivative of each element in $\theta$ is then the sum of partial derivatives over all outgoing edges from the corresponding \textit{$\theta$-node}.

The following briefly outlines a single Backpropagation step, i.e computing the partial derivative for a single edge $\textbf{u}\rightarrow\textbf{v}$ (notated as $\delta_{\textbf{u}\rightarrow\textbf{v}}$):   

\begin{equation}
    \delta_{\textbf{u}\rightarrow\textbf{v}} = \frac{\partial o_{\textbf{v}}}{\partial o_{\textbf{u}}} \cdot
    \begin{cases}
     \sum\limits_{e \in OutEdges(\textbf{v})} \delta_{e}     & \text{if $OutEdges(\textbf{v}) \neq \emptyset$}\\
    1       & \text{otherwise}
    \end{cases}
    \label{backprop}
\end{equation}

Where $OutEdges(\textbf{v})$ is the set out all outgoing edges from node $\textbf{v}$. The partial derivative of $t(\theta)$ w.r.t $\theta_i$, is then:

$$\frac{\partial t}{\partial \theta_i} = \sum\limits_{e \in OutEdges(\textbf{w})} \delta_{e} $$ 

 Where $\textbf{w}$ is the corresponding $\theta$\textit{-node} for $\theta_i$. Backpropagation thus can be used to compute the gradient for the loss of a given training sample in a feed-forward network, provided that the activation functions and loss function are easy to differentiate\footnote{This also applies for convex functions which are non-differentiable (e.g $ReLU$), in which case a subgradient is used.}. The process of computing the gradient on a 3-layered network is described in figure \ref{fig:backprop}.

\begin{figure}
    \centering
    \includegraphics{Figures/Electron.pdf}
    \caption{Backpropagation demonstrated on a 3-layered feed forward neural network}
    \label{fig:backprop}
\end{figure}

\subsection{Commonly Used Loss Functions}

A loss function $\mathcal{L}(\hat{y}, y)$ computes a (usually non-negative) value, measuring the inconsistency between a prediction $\hat{y}$ and ground truth $y$ for a given sample. 
As this work deals mostly with classification tasks, the following will describe two very common loss functions used for classification: the Cross Entropy loss, and the Hinge loss.

\subsubsection{Cross-Entropy Loss}
The Cross-Entropy loss function assumes the predicted vector $\hat{y}_{1xn}$ is a probability distribution, where $\hat{y}_i$ is the probability assigned to the $i$-th class from a set of $n$ classes. It is usually modeled by using Softmax as the output layer of the network. The ground truth $y$ is therefore a one-hot vector, with a probability of 1 assigned to the correct label (and 0 for the rest). The Cross-Entropy loss is the defined as the cross-entropy between the true distribution and the predicted one:

\begin{equation}
    \mathcal{L}_{CE}(\hat{y}, y) = \sum_{i=1}^{n} -y_{i}log(\hat{y}_i)
    \label{crossentropy}
\end{equation}

In the described setting, the Cross-Entropy boils down to simply taking the (minus) log of the probability assigned by the model to the correct label. It is worth noting, however, that Cross-Entropy loss can be used just as well for settings where the ground truth is not a single correct label, but a \emph{probability distribution} over the set of labels. 
A key property of Cross-Entropy loss is that it tends to infinity as the predicted probability tends to zero, rapidly penalizing the model for predictions that are both incorrect and confident.

\subsubsection{Hinge Loss}
The Hinge loss is suitable for scenarios where the model assigns each candidate with an unbounded score (rather than a probability), and the highest scoring candidate is considered as the predicted candidate. It takes a marginal approach, by penalizing predictions where the score assigned to the correct prediction is not larger then the rest by at least 1. The Hinge loss then corresponds to the score difference that violates this constraint mostly:

\begin{equation}
    \mathcal{L}_{Hinge}(\hat{y}, y) = max(0, 1 + max(\{\hat{y}_{i} | i \neq t\}) - \hat{y}_t)    \label{hingeloss}
\end{equation}

Where $t$ is the index of the correct candidate (obtained from $y$). Predictions assigning the correct candidate with the highest score, by a margin of at least 1, will receive a loss of 0. Otherwise, the loss is the violating difference:  $ 1 + \hat{y}_p - \hat{y}_t $ where $p$ is the index corresponding to the highest scoring candidate.

\subsection{Model Regularization and Overfitting}

As with any machine learning model, neural networks are also capable of overfitting the training data. This is especially true for deep networks, which have a very large amount of parameters and therefore more complexity. Simply picking the set of parameters that minimizes the loss on the training set is likely to overfit, at least to some degree. The common method of applying a regularization term (e.g $l1$, $l2$ regularization) in addition to the loss, can also be applied to the training process of a neural network. First, we define the \emph{cost} of a parameter vector $\theta$ w.r.t a of training sample $(\vec{x_i}, \vec{y_i})$ as the sum of its loss and a regularization term $R(\theta)$:

$$ Cost(\theta, \vec{x_i}, \vec{y_i}) = \mathcal{L}(\textbf{f}(\vec{x_i}, \theta), \vec{y_i}) + \lambda \cdot R(\theta) $$

Where $\lambda$ is a hyperparameter. Then, we modify the training process by minimizing the average cost instead of the average loss, so the function $t$ subject to minimization becomes:

$$ t(\theta) = \frac{1}{|Train|}\sum_{(\vec{x_i}, \vec{y_i}) \in Train} Cost(\theta, \vec{x_i}, \vec{y_i}) $$

Assuming that $R$ can be differentiated efficiently, the training process can just as well be applied to the modified $t$.

\subsubsection{Validation Loss and Early Stopping}

A common machine learning methodology for detection and prevention of overfitting is the use of a validation set. A validation set is a set of samples held out of the training set, on which the performance of a trained model is measured. As the validation samples had not been used for training, a high performance on the validation set is an indication of the model being able to generalize well. On the other hand, a model with a high performance on the training set and a low performance on the validation set, had likely to overfit. 

A validation set can be further utilized for prevention of overfitting, by examinig the loss on the validation set (the validation loss) throughout training. As training progresses, the loss on the training set generally decreases until reaching a local minima. The validation loss is also expected to decrease, until the model had overfitted the training set. At this point, the training loss may keep decreasing, but the validation loss may increase due to overfitting. A detection of such increase in validation loss may imply that the model had overfitted, and further training could hurt performance. Training is stopped at this point, and the set of parameters with the lowest validation loss seen throughout training is chosen. This adjustment to the training process is known as Early-Stopping. 

Detecting the point where the validation loss begins to increase due to overfitting is not so trivial. While the training loss is expected to decrease in general, specific training steps may increase it due to the fact the (a) the gradient of the loss is only estimated, and (b) the learning rate may be too big. Similarly, the validation loss may increase in some training steps before decreasing into its optimal value. Therefore, a single increase in validation loss is not enough. Some variations of early stopping does not allow stopping before a certain amount of epochs is completed. Others will stop only when the validation loss had not been decreased below its minimum in the last window of iterations. Another variation allows the training process to continue until the training loss stabilizes, and only then the best parameters are chosen. This variation have no risk of stopping too early, at the cost of a longer training process.

% \begin{tikzpicture}
% \begin{axis}[
%     axis lines = left,
%     xlabel = $x$,
%     ylabel = {$f(x)$},
% ]
% %Below the red parabola is defined
% \addplot [
%     domain=-10:10, 
%     samples=100, 
%     color=red,
% ]
% {x^2 - 2*x - 1};
% \addlegendentry{$x^2 - 2x - 1$}
% %Here the blue parabloa is defined
% \addplot [
%     domain=-10:10, 
%     samples=100, 
%     color=blue,
%     ]
%     {x^2 + 2*x + 1};
% \addlegendentry{$x^2 + 2x + 1$}
 
% \end{axis}
% \end{tikzpicture}



\subsubsection{*Neural nets overview}
\subsubsection{*Activation functions}
\subsubsection{*Loss functions}
\subsubsection{*Training}
\subsubsection{*Hyperparameters vs. Model parameters, tuning}
\subsubsection{*Regularization and drop-out}
\subsection{Feed Forward Networks}
\subsection{Recurrent Neural Networks}
\subsubsection{Bidirectional RNN}
\subsubsection{Stacked RNNs}
\subsubsection{Long-Short-Term-Memory (LSTM)}
\subsection{Lookup Tables}
\subsection{The Stacked BiLSTM-MLP-Softmax Architecture}
\subsection{Neural Networks in this work}
\pagebreak

%-----------------------------------
%	SECTION 3
%-----------------------------------
\section{Word Embeddings Overview}
\subsection{Main Approches}
\subsubsection{The Skip-Gram Model (Word2Vec)}
\subsubsection{Global Vectors for Word Representation (GLoVe)}
\subsubsection{Syntactic Embeddings}
\subsubsection{FastText?}
\subsection{Contextualized Embeddings}
\subsubsection{ELMo}
\subsection{Multi-Lingual Word Embeddings}
\subsubsection{MUSE}
\subsection{Learning via Word Embeddings}
\subsubsection{*Basic usage with a FF-Network}
\subsubsection{*Updating embeddings through training}
\subsubsection{*Contextutalized representations with RNNs}
\subsection{Handling Out-Of-Vocabulary words}
\subsubsection{*By stemming}
\subsubsection{*By learning an internal representation}
\subsubsection{*Character-based embeddings}
\subsection{Use of Word Embeddings in this work}
\pagebreak


%-----------------------------------
%	SECTION 4
%-----------------------------------
\section{Natural Language Parsing and the Prepositional Phrase Attachment Problem}
\subsubsection{*Outline}
\subsection{Sentence Parsing}
\subsubsection{Dependency vs. Constituency Based Parsing}
\subsubsection{Syntactic vs. Semantic Parsers}
\subsection{The PP-Attachment problem}
\subsubsection{*Overview, examples}
\subsection{PP Attachment Disambiguation Approaches}
\subsubsection{As a sub-task of sentence parsing}
\subsubsection{As an isolated task}
\subsection{PP Attachment and PSS}
\subsubsection{*Motivation of using PSS for PPA}
\pagebreak

%-----------------------------------
%	SECTION 5
%-----------------------------------

%-----------------------------------
%	SECTION 6
%-----------------------------------
\section{Universal Dependencies}
\pagebreak

%-----------------------------------
%	SECTION 7
%-----------------------------------
\section{Named Entity Recognition}
